\twocolumn[  
\begin{@twocolumnfalse}
    \begin{center}
        \noindent
        \begin{minipage}{1.25\columnwidth}
        \setlength{\parindent}{15pt}
        \mbox{}\vfill
        \section*{Preface}

The \CENTREX\ experiment aims to provide a 30-fold increase in sensitivity over the best present upper bounds on the strength of time-violating fundamental interactions by measuring shifts in nuclear resonant frequencies in thallium fluoride. In this prospectus, I describe the motivation for studying these time-violating interactions, and provide an overview of past work. I present the current measurement scheme for this experiment. Since the apparatus is still under construction, I present the components already built, and estimate the timeline for completion of the experiment.

Most of the work will be carried out at the Sloane Physics Laboratory, in the spaces of Prof DeMille. Parts of the experiment have been, and will be, constructed at UMass and Columbia --- for example, parts of the laser systems, and the interaction region electrodes and vacuum chamber.

For my thesis, I expect to work to some extent on all parts of the apparatus, but especially the following:
\begin{itemize}
    \item beam source (which I have already built myself); see pages \pageref{sect:beamsources}ff and \pageref{sect:beamsource}ff
    \item rotational cooling and other microwave systems (I have already assembled the prototypes of the systems, and partially they are already used in the preliminary testing); see pages \pageref{sect:rotcool}ff
    \item magnetic shielding (which will be mine to characterize and assemble --- there is only a prototype so far)
    \item simulations of non-adiabatic transitions (pages \pageref{sect:nonadiabatic}ff)
    \item simulating rotational cooling or other aspects of the experiment, as needed (TBD).
\end{itemize}
%
Rough timeline to finishing thesis at end of my 6th year:
%
\begin{itemize}
  \item 2020: Finish characterization of the rotational cooling. Start construction state-preparation regions. Start construction of optical cycling and detection.
  \item 2021: Testing of state-prep, optical cycling, detection regions. Testing interaction region, and final magnetic shielding.
  \item 2022: First experimental runs; taking data; data analysis; error characterisation.
  \item 2023: Getting to a first measurement result.
\end{itemize}
%
            Please refer to the table of contents, following page, for the organization of the rest of this document. The thesis is expected to have a similar structure, plus chapters on data analysis, interpretation, and conclusion. Some sections are empty or are marked incomplete, indicating the work that has not yet been performed for the thesis.

        \end{minipage}
        \vfill
    \end{center}
    \vspace{1cm}
\end{@twocolumnfalse}
]
